\documentclass[12pt]{article}
\usepackage[utf8]{inputenc}
\usepackage{amsmath, amssymb, amsthm}
\usepackage{geometry}
\usepackage{fancyhdr}

\geometry{letterpaper, margin=1in}
\setlength{\headheight}{15pt}

% Header setup
\pagestyle{fancy}
\fancyhf{}
\lhead{Math 477 Spring 2026}
\rhead{Jeffrey Lu}
\cfoot{\thepage}

\title{Homework Assignment 1 \\ \large Math 292, Spring 2026 \\ Rutgers University}
\author{Jeffrey Lu}

\begin{document}

\maketitle

\noindent \textbf{Problem 1:}
\begin{enumerate}
    \item[(a)]
          Let sample space be defined by
          \[
              S = \{1, 2, 3, 4 5, 6\}
          \]
    \item[(b)]
          \[
              A\cap{B} = \{1,3 \}
          \]
          \[
              (A\cup{B})^c = \{4,6\}
          \]
    \item[(c)]
          \begin{align*}
              P(A\cap{B})     & = \frac{|A\cap{B}|}{|S|} = \frac{1}{3}     \\
              P((A\cup{B})^c) & = \frac{|(A\cup{B})^c|}{|S|} = \frac{1}{3}
          \end{align*}
\end{enumerate}

\noindent \textbf{Problem 2:}
Using the multinomial coefficient, we get
\begin{align*}
    \binom{11}{2~4~4} = \frac{11!}{4!*4!*2!} & = 34,650 \text{ different letter arrangements}
\end{align*}

\noindent \textbf{Problem 3:}
\begin{enumerate}
    \item[(a)]
          \[
              \binom{8}{5} - \binom{6}{3} = 56 - 20 = 36 \text{ choices}
          \]
    \item[(b)]
          \[
              \binom{6}{5} + \binom{6}{3} = 6 + 20 = 26 \text{ choices}
          \]

\end{enumerate}

\noindent \textbf{Problem 4:}
\begin{enumerate}
    \item[(a)] Using stars and bars we get
          \[
              \binom{13}{4} = 715 \text{ possible divisions}
          \]

    \item[(b)] We subtract 5 from n since each school must receive one and there is 5 schools
          \[
              \binom{8}{4} = 70 \text{ possible divisions}
          \]
\end{enumerate}

\noindent \textbf{Problem 5:}
\begin{enumerate}
    \item[(a)]
          \[
              15^3 = 3375 \text{ different cones}
          \]
    \item[(b)]
          \[
              P(15,3) = 2730 \text{ different cones}
          \]
    \item[(c)]
          \[
              \binom{17}{3} = 680 \text{ different cones}
          \]
    \item[(d)]
          \[
              \binom{15}{3} = 445 \text{ different cones}
          \]
\end{enumerate}

\noindent \textbf{Problem 6:}
\begin{enumerate}
    \item[(a)]
          \[
              8! = 40320 \text{ different seating arrangements}
          \]
    \item[(b)]
          \[
              2! * 7! =  100800 \text{ different seating arrangements}
          \]
    \item[(c)]
          \[
              2! * 4! * 4! =  1152 \text{ different seating arrangements}
          \]
    \item[(d)]
          \[
              4! * 5! = 2880 \text{ different seating arrangements}
          \]
    \item[(e)]
          \[
              4! * (2!)^4 =384 \text{ different seating arrangements}
          \]
\end{enumerate}
\end{document}